% !TEX root = ../main.tex
\section{Limitations and conclusion}
\label{sec:09_conclusion}

The value function of a reward-guided diffusion is a conditional log-partition
function, and the only unbiased Monte-Carlo targets for it live in exponential
space. That fact is what forces the awkward regression, and squared error
handles it badly in a specific and asymmetric way: over-predictions explode,
which clipping survives, and under-predictions vanish, which nothing survives.
Since every Bregman divergence has the value as its minimizer, the potential is
free, and in log space the entire family reduces to one positive weight on the
residual. Choosing that weight to be $\ln y/(y-1)$ gives gradients linear in the
prediction on both sides, in closed form, with a numerically stable
implementation. It is a drop-in replacement for a loss function.

What the measurements support is narrower than that construction might suggest,
and we would rather state it than let the framing carry it. The Spence weight
repairs exponential-space squared error on both metrics at every dimension we
ran, and beats generalised KL on value fidelity, so controlling both tails is
demonstrably better than controlling one. It does not dominate the family. We ran the
experiment that would decide this --- the same grid with gradient clipping on
every arm --- and it decided against us: clipped Itakura--Saito is the better
controller above $d=2$ and ties on value, and clipping turns out to help our
loss substantially too, which a weight that genuinely needed no clipping should
not do. The durable contribution is therefore the reduction itself --- that
loss design here is the choice of a single positive weight, and that squared
error is a poor choice of it --- and not a claim that our weight is the best
available one. We would recommend the reduction to a practitioner, and clipped
Itakura--Saito as the loss.

\paragraph{Limitations.}
The empirical claims are narrower than the analysis.

\emph{Reward-call budget.} All our targets have cheap analytic rewards, so we
measure cost in samples seen. Under an expensive reward --- a quantum-chemistry
calculation, a docking score, a simulator rollout --- the binding budget is
total evaluations of $r$, and the ranking between off-policy anchors and
on-policy sampling could differ from what we report
(\S\ref{subsec:targets}).

\emph{Transfer.} The on-policy recipe was fitted on one benchmark family, and
does not transfer unchanged to the Boltzmann targets of
\S\ref{subsec:gmm40}. We report the fitted constants as an existence proof that
a single family can cover the dimension range, not as values anyone should
reuse.

\emph{Many-Well-32.} We do not reach a competitive number on this target and
have diagnosed but not fixed the cause (\S\ref{subsec:gmm40}). Until that is
closed, the GMM-40 result should be read as one favourable external
target rather than as a general claim about Boltzmann sampling.

\emph{No pretrained model.} This is the largest gap. \S\ref{sec:02_setting}
observes that if the base process is a pretrained diffusion or flow model then
$f$ is its learned drift and nothing in the development changes, and we believe
that, but we never test it: every experiment here uses analytic base processes
--- Gaussian-mixture interpolants and driftless Brownian motion --- with
residual MLPs at gradient batch size $4$ on synthetic tilts. Whether the
conditioning differences we measure survive at the scale and with the learned
drifts of an actual reward-guided diffusion is untested, and a reader should
weight the results accordingly.

\emph{Scope of the benchmark.} The dimensional-scaling study fixes headroom at
$6$ nats and caps network width at $256$ for $d\ge8$. Both the crossover
dimension between on- and off-policy training and the size of the loss effect
plausibly move with these, and we have not varied them.

\emph{What the loss does not do.} The Spence loss improves conditioning. It
does not make rare high-value targets easier to find, which is the separate
problem of \S\ref{sec:06_sparsity_onpolicy} and the reason the two halves of
this paper are both necessary.

\paragraph{Outlook.}
The gradient-weight view suggests its own next question. Once the design space
is a single positive function $w$ on $(0,\infty)$, the weight can be chosen for
things other than tail behaviour --- robustness to target noise, or a
prescribed influence function --- and \eqref{eq:req-over}--\eqref{eq:req-under}
pin down only the asymptotics, leaving a family of admissible weights that we
have not explored.
